% interaction/composition.tex — append, liftAppend, replicate, stateChain, Chain

\section{Sequential Composition}\label{sec:interaction-composition}

Sequential composition is the primary way protocols grow: run one interaction,
then continue with a second whose shape may depend on the outcome of the first.
The key technical tool in this section is the $\mathsf{liftAppend}$
combinator, which eliminates a cascade of type casts that would otherwise infect
every downstream definition.

\subsection{Dependent append}

\begin{definition}[Spec.append]
  \label{int:spec-append}
  Given $s_1 : \mathsf{Spec}$ and $s_2 : \Transcript\;s_1 \to \mathsf{Spec}$,
  the \emph{dependent append} $s_1.\mathsf{append}\;s_2$ fuses the two
  interactions into one.  Informally, at each leaf of~$s_1$ the corresponding
  $s_2\;\mathit{tr}_1$ is grafted on.
  \lean{Interaction.Spec.append}
  \uses{int:spec}
\end{definition}

\begin{definition}[Transcript.append / split]
  \label{int:transcript-append-split}
  $\Transcript.\mathsf{append}$ combines a first-phase transcript $\mathit{tr}_1$
  and a second-phase transcript $\mathit{tr}_2$ into a transcript of the composed
  interaction.  $\Transcript.\mathsf{split}$ is its inverse: it decomposes a
  transcript of $s_1.\mathsf{append}\;s_2$ into its two components.  The
  round-trip lemmas $\mathsf{split\_append}$ and $\mathsf{append\_split}$ hold.
  \lean{Interaction.Spec.Transcript.append}
  \lean{Interaction.Spec.Transcript.split}
  \uses{int:spec-append, int:transcript}
\end{definition}

\subsection{The \texorpdfstring{$\mathsf{liftAppend}$}{liftAppend} combinator}

When composing strategies, the output type of the second phase depends on the
first-phase transcript.  A natural formulation is a two-argument family
$F : \Transcript\;s_1 \to \Transcript\;(s_2\;\mathit{tr}_1) \to \Type$.
But strategies on $s_1.\mathsf{append}\;s_2$ need a \emph{single-argument}
family on $\Transcript\;(s_1.\mathsf{append}\;s_2)$.

\begin{definition}[Transcript.liftAppend]
  \label{int:lift-append}
  $\mathsf{liftAppend}\;s_1\;s_2\;F$ lifts $F$ to a single-argument family on
  $\Transcript\;(s_1.\mathsf{append}\;s_2)$.  The crucial property is
  \emph{definitional computation}:
  \[
    \mathsf{liftAppend}\;s_1\;s_2\;F\;
      (\Transcript.\mathsf{append}\;s_1\;s_2\;\mathit{tr}_1\;\mathit{tr}_2)
    \;\equiv\;
    F\;\mathit{tr}_1\;\mathit{tr}_2
  \]
  where $\equiv$ denotes judgmental equality---no explicit cast or transport is
  needed.
  \lean{Interaction.Spec.Transcript.liftAppend}
  \uses{int:spec-append, int:transcript}
\end{definition}

This property propagates through the entire stack. $\mathsf{stateChainFamily}$
(Section~\ref{sec:iteration}) uses $\mathsf{liftAppend}$ at each stage of a
state chain; $\mathsf{Chain.outputFamily}$ uses it at each round of a
continuation chain; and all strategy composition combinators and security
composition theorems factor through it.  Without $\mathsf{liftAppend}$, every
such combinator would require explicit casts between the two-argument and
single-argument views.

Companion operations include $\mathsf{packAppend}$ and
$\mathsf{unpackAppend}$ (transport between the two views),
$\mathsf{liftAppendRel}$ and $\mathsf{liftAppendPred}$ (lift binary
relations and predicates for use in security statements).

\subsection{Strategy composition}

\begin{definition}[Strategy.comp]
  \label{int:strategy-comp}
  $\Strategy.\mathsf{comp}\;s_1\;s_2$ composes two strategies along
  $\mathsf{Spec.append}$.  The continuation receives the first phase's output
  and produces a second-phase strategy.  The composed output type is given by
  $\mathsf{liftAppend}$:
  \[
    \Strategy.\mathsf{comp} :
      \Strategy\;m\;s_1\;\mathit{Mid} \to
      \bigl(\forall\;\mathit{tr}_1,\;
        \mathit{Mid}\;\mathit{tr}_1 \to
        m\,(\Strategy\;m\;(s_2\;\mathit{tr}_1)\;(F\;\mathit{tr}_1))\bigr)
      \to m\,\bigl(\Strategy\;m\;(s_1.\mathsf{append}\;s_2)\;
        (\mathsf{liftAppend}\;s_1\;s_2\;F)\bigr).
  \]
  A ``flat'' variant $\mathsf{compFlat}$ uses a single output family directly.
  $\mathsf{splitPrefix}$ decomposes a composed strategy into its prefix and
  continuation.
  \lean{Interaction.Spec.Strategy.comp}
  \uses{int:spec-append, int:strategy, int:lift-append}
\end{definition}

\subsection{Iteration mechanisms}\label{sec:iteration}

Protocols frequently iterate a single-round interaction multiple times.
We provide three iteration mechanisms at increasing levels of generality.

\begin{definition}[Spec.replicate]
  \label{int:spec-replicate}
  $\mathit{spec}.\mathsf{replicate}\;n$ is the $n$-fold non-dependent append of
  the same $\mathsf{Spec}$.  The continuation ignores the transcript of each
  round.
  \lean{Interaction.Spec.replicate}
  \uses{int:spec-append}
\end{definition}

\begin{definition}[Spec.stateChain]
  \label{int:spec-state-chain}
  $\mathsf{stateChain}\;\mathit{Stage}\;\mathit{spec}\;\mathit{advance}\;n\;i\;s$
  is an $n$-stage state-indexed composition.  At each stage~$i$ with
  state $s : \mathit{Stage}\;i$, the interaction is
  $\mathit{spec}\;i\;s$; after the stage completes with
  transcript~$\mathit{tr}$, the state advances to
  $\mathit{advance}\;i\;s\;\mathit{tr} : \mathit{Stage}\;(i+1)$.
  \lean{Interaction.Spec.stateChain}
  \uses{int:spec-append}
\end{definition}

\begin{definition}[Spec.Chain]
  \label{int:spec-chain}
  A \emph{chain} $\mathsf{Chain}\;n$ is a depth-indexed telescope: at each
  level it carries the current round's $\mathsf{Spec}$ and, for each possible
  transcript, the recipe for the remaining rounds.  There is no external state
  type.
  \[
    \mathsf{Chain}\;0 = \Unit, \qquad
    \mathsf{Chain}\;(n+1) =
      (\mathit{spec} : \mathsf{Spec}) \times
      (\Transcript\;\mathit{spec} \to \mathsf{Chain}\;n).
  \]
  $\mathsf{Chain.toSpec}$ converts a chain to a concrete $\mathsf{Spec}$ via
  iterated $\mathsf{append}$.
  \lean{Interaction.Spec.Chain}
  \uses{int:spec, int:spec-append}
\end{definition}

These three mechanisms are related by specialization:
%
\begin{center}
\begin{tabular}{l c c l}
  \textbf{Mechanism} & \textbf{State?} & \textbf{Transcript-dep.?} & \textbf{Primary use} \\
  \hline
  $\mathsf{replicate}$ & No & No & Uniform rounds \\
  $\mathsf{stateChain}$ & Yes ($\mathit{Stage}\;i$) & Yes & State machines \\
  $\mathsf{Chain}$ & No (baked in) & Yes & Continuation-style \\
\end{tabular}
\end{center}

$\mathsf{Chain}$ is the most fundamental: $\mathsf{Chain.replicate}$ recovers
$\mathsf{Spec.replicate}$, and $\mathsf{Chain.ofStateMachine}$ recovers
$\mathsf{Spec.stateChain}$.

\begin{remark}
  The $\mathsf{Chain}$ examples in the formalization include a
  \emph{growing-messages} protocol (round~$k$ exchanges a value from
  $\Fin\;(k+1)$) and a \emph{prefix-dependent} protocol where the third-round
  message type genuinely depends on both prior moves.  These demonstrate that
  transcript dependence is expressible without any external state.
\end{remark}
